\documentclass[10pt,twocolumn,letterpaper]{article}

\usepackage{graphicx} 

\usepackage[pagebackref=true,breaklinks=true,letterpaper=true,colorlinks,bookmarks=false]{hyperref}

\begin{document}
    \title{Project Report : CS 7643}

    \maketitle

    \begin{abstract}
        Abstract
    \end{abstract}

    \section{Introduction/Motivation}
    The incidence of skin cancer in the U.S. is quite significant, more than 
    9,500 people are diagnosed with skin cancer every day, making it the most 
    common form of cancer in the country. It's estimated that one in five 
    Americans will develop skin cancer by the age of 70. Notably, higher 
    mortality rates are observed among men, non-Hispanic Whites, and those 
    aged over 70 \cite{Authors1}. \\
    \\
    Although various deep learning approaches have been used for 
    computer-based skin cancer detection in recent years. Most detection
    algorithms use convuolution nerual networks (CNN) to
    classify skin leasions.  To the best of our knowledge, few deep learning 
    studies on skin cancer have integrated metadata such as age, sex or leasion locations; 
    only one study has utilized the dataset incorporating the additional risk 
    factor such as race (the Caucasian and xanthous races dataset) to include these 
    meta-features \cite{Authors2, Authors3, Authors4}.\\
    \\
    In our research, we aim to incorporate these clinically investigated 
    risk factors into an enhanced model scheme \cite{Authors5}. By applying CNN transfer 
    learning to the classification task along with these additional risk 
    factors, we seek to refine the existing model architecture. This approach will improve the differentiation of malignant and 
    benign lesions, boost early detection of skin cancer, enhance diagnostic 
    accuracy, and ultimately improve clinical outcomes. \\
    \\
    Improving accuracies and early detections would be important for doctors
    patients, hospitals, and insurance companies.  Early detection can give the
    patient a 99\% survival rate \cite{Authors6}.  This also means that early
    detection results in less hospital stays, smaller operations, and less
    expenditures by patients and insurance.  Having such tools also aids
    doctors by giving them another tool to determine if a leasion is serious
    or not.  Furthermore, better accuracies will not only catch potentially malignant
    leasions but also determine if a leasion is begin and does not need to
    be removed. \\
    \\
    This research will be utilizing the HAM10000 ("Human Against Machine 10000 images")
    dataset and the ISIC 2024 dataset, both of which can be found on kaggle.com.  These datasets come with one to many images
    of leasions that can be used for training machine learning models along with
    metadata about the patient to accompany these images.  The metadata that accompanies
    the images will be the primary focus of this research as this will add new features
    in an attempt to increase accuracy. \\
    
    \section{Approach}
    Background research identified several pretrained models that were shown to perform
    well on the HAM10000 data set \cite{Authors7}.  That goal of this research is to incorporate patient
    metadata with the models to see how it improved on the accuracy, thus, the ResNet101 model
    that is in pytorch was selected.  The hypothesis is that adding additional data about the
    patients age, sex, and leasion locations would increase the accuracies.  For example, if you
    look at Figure~\ref{fig:agebox}, it can be seen that the risk for malignant increase with age.
    The final fully connected layer that did classification was
    changed to take in the metadata along with the results from the ResNet101 model. See \\
    \begin{figure}[t]
        \begin{center}
           \includegraphics[width=0.8\linewidth]{plots/agebox.png}
        \end{center}
        \caption{Benign and malignant cases by age.}
        \label{fig:agebox}
    \end{figure}
    A few problems came up at this point of the research.  First, many of the leasions
    contained multiple images, at different zoom levels, slightly rotated, and at different angles.
    The first iterations of our model we filtered out these images and choose only one images for
    each leasion.  The worry was that the multiple images would bias add bias.  The next problem
    with the data is that this is a very unbalanced dataset.  There are significantly more cases
    of begin then malignant photos as shown in Figure~\ref{fig:distbysex}.  To deal with the
    unbalanced dataset different loss functions were explored both cross-entropy and focal loss
    functions. \\
    \begin{figure}[t]
        \begin{center}
           \includegraphics[width=0.8\linewidth]{plots/distbysex.png}
        \end{center}
        \caption{Distribution of benign and malignant cases by sex.}
        \label{fig:distbysex}
    \end{figure}
    Two models were created for this research.  In our background research most pretrained models had pretty similar
    results and ResNet101 was selected to be our benchmark for comparison \cite{Authors7,Authors8,Authors9,Authors10,Authors11,Authors12}.
    Next, the same model that was used as a benchmark was used again and implemented to incorporate
    the additional parient metadata (age, sec, leasion location) as additional input features.  The Figure~\ref{fig:arch} below
    shows the architecture of the model with the addition of patient metadata. Finally, we then attempted to apply a
    second data set to our model to see how it performed in transfer learning. \\
    \begin{figure}[t]
        \begin{center}
           \includegraphics[width=0.8\linewidth]{plots/arch.png}
        \end{center}
        \caption{Model architecture with patient metadata added.}
        \label{fig:arch}
    \end{figure}

    \section{Approach}



    {\small
        \bibliographystyle{ieee_fullname}
        \bibliography{report}
    }

\end{document}